% Figure: a design pseudo-acceleration response spectrum, the peak
% acceleration of a SDOF oscillator versus its natural period for a
% fixed damping ratio. Three regions: rising at short period, a flat
% acceleration-controlled plateau, and a falling velocity/displacement
% branch. Built from a piecewise expression so it needs no data file.
\begin{figure}[htbp]
\centering
\begin{tikzpicture}
\begin{axis}[
  width=12cm, height=7.2cm,
  xlabel={oscillator natural period $T$ (s)},
  ylabel={pseudo-acceleration $S_a$ ($g$)},
  xmin=0, xmax=4, ymin=0, ymax=1.05,
  axis lines=left,
  samples=400,
  every axis x label/.style={at={(current axis.right of origin)}, anchor=west},
  every axis y label/.style={at={(current axis.above origin)}, anchor=south},
  grid=major, grid style={enggray!20},
  tick label style={font=\scriptsize},
  label style={font=\small},
  legend style={font=\scriptsize},
  legend pos=north east,
]
\addplot[engblue, very thick, domain=0:0.15, samples=2] {0.4 + (0.9-0.4)/0.15*x};
\addplot[engblue, very thick, domain=0.15:0.5, samples=2, forget plot] {0.9};
\addplot[engblue, very thick, domain=0.5:4, samples=200, forget plot] {0.45/x};
\addlegendentry{$\zeta=0.05$ design spectrum}
\addplot[engred, dotted, thick, forget plot] coordinates {(2.0,0) (2.0,0.225)};
\addplot[engred, only marks, mark=*, mark size=2pt, forget plot] coordinates {(2.0,0.225)};
\node[engred, font=\scriptsize, anchor=west] at (axis cs:2.05,0.30) {$T_1=2$ s: $S_a=0.225\,g$};
\node[enggray, font=\scriptsize, anchor=south] at (axis cs:0.32,0.92) {acceleration plateau};
\node[enggray, font=\scriptsize, anchor=west] at (axis cs:1.4,0.55) {velocity branch};
\end{axis}
\end{tikzpicture}
\caption[A design response spectrum]{A design pseudo-acceleration
response spectrum: the peak acceleration of a single-degree-of-freedom
oscillator of five-percent damping, plotted against its natural period,
for one design ground motion. A short-period (stiff) structure rides
the ground almost rigidly and sees the plateau acceleration; a
long-period (flexible) structure filters the high-frequency ground
content and sees less, the falling velocity branch where
$S_a\propto 1/T$. The spectrum collapses an entire time-history
analysis into one lookup: read the structure's fundamental period off
the horizontal axis, read the pseudo-acceleration off the curve,
multiply by the participating modal mass to get the base shear. The
example structure at $T_1=2\,\si{\second}$ draws $0.225\,g$, a quarter
of what a stiff structure on the same site would feel, which is why
long-period framing is one route to lower seismic demand.}
\label{fig:vol03ch06:response-spectrum}
\end{figure}
